\documentclass[a4paper, 12pt]{article}
\usepackage[utf8]{inputenc}
\usepackage[spanish, provide=*]{babel}
\usepackage{amsmath, amssymb}
\usepackage{graphicx}
\usepackage{hyperref}
\usepackage{listings}
\usepackage{tcolorbox}
\usepackage{xcolor}
\usepackage{fancyhdr}
\usepackage{geometry}

\geometry{left=2.5cm, right=2.5cm, top=3cm, bottom=4.5cm}

% Para las tablas
\usepackage{float}  %Para que no se cambien de lugar
\usepackage{pgfplotstable}
\usepackage{array}
\usepackage{colortbl}
\usepackage{url} % para que el campo note muestre URLs bien

% Estilo
\usepackage{lastpage}

% Definición de colores
\definecolor{cafetitulos}{RGB}{150,75,0}
\definecolor{redv}{RGB}{165, 43, 42}
\definecolor{cafe}{RGB}{130,75,0}
\definecolor{cafet}{RGB}{204, 121, 5}

% Configuración de encabezados y pies de página
\pagestyle{fancy}
\renewcommand{\headrulewidth}{0pt} % Elimina la línea por defecto del encabezado

% Encabezado
\fancyhead[L]{\includegraphics[height=1.5cm]{logomanual.png}}
\fancyhead[R]{\nouppercase{\leftmark}} % Nombre de la sección

% Pie de página
\renewcommand{\footrulewidth}{0.5pt} % Línea separadora
\renewcommand{\footrule}{\hbox to\headwidth{\color{cafe}\leaders\hrule height \footrulewidth\hfill}}
\fancyfoot[L]{Documento de Diseño y Arquitectura - 2025}
\fancyfoot[C]{}
\fancyfoot[R]{Página \thepage} % Número de página a la derecha

% Ajustar altura del encabezado y espacio para el pie
\setlength{\headheight}{52pt} % Altura del encabezado (ajusta según el logo)
\setlength{\footskip}{40pt} % Espacio entre el texto y el pie (evita solapamiento)

% Titulos
\usepackage{titlesec} % Para personalizar section/subsection
\titleformat{\section}
  {\normalfont\Large\bfseries}{\thesection}{1em}{\color{cafetitulos}} % Section en café
\titleformat{\subsection}
  {\normalfont\large\bfseries}{\thesubsection}{1em}{\color{cafe}}
\titleformat{\subsubsection}{\normalfont\bfseries}{\thesubsubsection}{1em}{\color{cafetitulos}}

% Configuración para código fuente
\lstset{
    language=Java, % puedes cambiar a Python, SQL, etc.
    basicstyle=\ttfamily\small,
    backgroundcolor=\color{gray!10},
    frame=single,
    breaklines=true,
    keywordstyle=\color{cafetitulos},
    commentstyle=\color{gray},
    stringstyle=\color{redv}
}


\title{
\includegraphics[width=5cm]{logo.png}\\

{\Large CHIBCHAWEB HOSTING PLATFORM}\\
{\Large APLICACIÓN PARA UNA COMPAÑÍA DE HOSPEDAJE DE PÁGINAS WEB}\\[0.5em]

\textbf{DOCUMENTO DE DISEÑO Y ARQUITECTURA}\\[0.5em]
}
\author{
\textbf{Grupo A}\\[0.5em]
Ingenieros:\\
Andrés José Acevedo Ardila\\
Brayan Steven Sánchez Casilima\\
Christian Camilo Lancheros Sánchez\\
Juan Esteban Ordoñez Sandoval\\
Karen Tatiana Bravo Rodríguez\\
Luis Felipe Mayorga Tibaquicha\\\\\\
Universidad Distrital Francisco José de Caldas\\
Facultad de Ingeniería\\
Ingeniería de sistemas\\
}
\date{Bogotá D.C., Agosto 2025}


% Inicio Documento
\begin{document}
\maketitle

\newpage
\begin{titlepage}
\newgeometry{left=2cm, right=0cm, top=2cm, bottom=2cm} % sin margen derecho
\thispagestyle{empty}

\begin{tikzpicture}[remember picture, overlay]
    \node[anchor=north east, inner sep=0pt, xshift=12cm] at (current page.north east) {
        \includegraphics[height=\paperheight]{registro-imagen.png}
    };
\end{tikzpicture}

\vfill
{\includegraphics[width=0.6\textwidth]{logomanual.png}\par}
\vfill
{\Huge Grupo A\par}
{\Huge Manual Técnico ChibchaWeb\par}
{\Large 2025\par}
\vfill


\end{titlepage}


\newpage
\phantomsection
\tableofcontents

\newpage
\section{Listado de Figuras}
\listoffigures

\newpage
\section{Listado de Tablas}
\listoftables

\newpage
\section{Descripción del Documento}
\subsection{Propósito y Audiencia}
Este documento tiene como propósito describir la arquitectura de software propuesta para el desarrollo de una solución integral que apoye los procesos operativos de ChibchaWeb, una compañía dedicada al hospedaje de páginas web con sede en Sugamuxi y proyección de expansión internacional.\\
El documento presenta los componentes clave del sistema, los requerimientos funcionales y no funcionales, los patrones de diseño seleccionados, y los modelos arquitectónicos que sustentan la solución. También detalla la metodología de desarrollo adoptada, las decisiones técnicas tomadas y su justificación.\\
Este documento está dirigido a los siguientes perfiles:

\begin{itemize}
    \item{Equipo de Desarrollo: Para comprender la arquitectura, los patrones y las tecnologías definidas que guiarán la implementación del sistema.}
    \item{Clientes y Stakeholders de ChibchaWeb: Para validar que los requerimientos funcionales y no funcionales del negocio han sido entendidos y serán satisfechos por el sistema.}
    \item{Equipo de QA y Soporte Técnico: Para conocer los aspectos técnicos relevantes de la solución y planificar las estrategias de prueba y soporte.}
    \item{Líderes de Proyecto y Gerencia: Para tomar decisiones informadas sobre tiempos, recursos y alineación con objetivos estratégicos.}
\end{itemize}

\subsection{Organización del Documento}
Este documento empieza con un panorama general del proyecto consistente en objetivos generales y específicos, seguidos de una contextualización del entorno y conceptos importante del proyecto, esto con el fin de generar un marco teórico que aclare y mejore la organización empresarial.\\
Luego de esto, se define el modelo de desarrollo y por lo tanto la organización del equipo de desarrollo. Después se tiene en cuenta las motivaciones y los riesgos relacionados con el desarrollo del proyecto, con el fin de estimar la utilidad. Además de las restricciones (tecnológico, negocio) y la gestión de calidad (con parámetros medibles para cumplir objetivos).\\
Por último, se definen los diagramas principales para empezar el proceso del producto, describiendo los diferentes modelos y entidades del desarrollo.
Para terminar, se hace una pequeña retroalimentación con el fin de reconocer errores dentro del desarrollo del proyecto y corregirlos en próximos proyectos.
\subsection{Convenciones}

\subsubsection*{Identificadores de requerimientos:}
\begin{itemize}
	\item Requerimientos funcionales: RF-XX (ej. RF-01).
\end{itemize}


\subsection{Terminología y Definiciones}

\begin{itemize}
\item Cliente: Usuario registrado que adquiere servicios de hosting o dominios.
\item Distribuidor: Usuario que intermedia la venta de dominios y recibe comisiones.
\item Empleado: Usuario interno con funciones operativas dentro de ChibchaWeb.
\item Administrador: Usuario con permisos totales de gestión.
\item Ticket: Registro de solicitud de soporte técnico.
\item Dominio: Nombre único utilizado para identificar un sitio web en Internet.
\item Carrito: Módulo de compra donde se almacenan temporalmente productos antes del pago.
\item Scrum: Marco de trabajo ágil usado para el desarrollo del proyecto.
\item Frontend: Parte de la aplicación con la que interactúa el usuario final.
\item Backend: Parte de la aplicación que maneja la lógica de negocio y la comunicación con la base de datos.
\item PSE: Plataforma de pagos electrónicos utilizada en Colombia.
\item XML: Lenguaje de marcado utilizado para intercambio de información estructurada.
\item Railway: Plataforma de despliegue utilizada para alojar el sistema.
\item Scraping: Técnica para extraer información de páginas web de terceros.
\item Método de pago: Forma registrada por el usuario para realizar transacciones (tarjeta de crédito, PSE, etc.).
\end{itemize}


\newpage
\section{Generalidades del proyecto}
\subsection{Problema por resolver}
ChibchaWeb, una empresa de hospedaje de sitios web ubicada en Sugamuxi, ha experimentado un crecimiento constante en su base de clientes, principalmente en Colombia y países cercanos, y se proyecta una expansión hacia otros continentes como África. Actualmente administra miles de sitios web bajo distintos paquetes de hosting (Chibcha-Platino, Chibcha-Plata y Chibcha-Oro) sobre plataformas Windows y Unix. Los clientes pueden escoger planes de pago flexibles (mensual, trimestral, semestral o anual), y aunque el método principal de pago es con tarjeta de crédito, se planea habilitar próximamente otros medios como PSE.

El modelo de negocio incluye además una red de distribuidores categorizados como BÁSICOS o PREMIUM, con diferentes esquemas de comisiones. ChibchaWeb también intermedia en el registro y transferencia de dominios a través de terceros registradores, ofreciendo este servicio adicional a sus clientes. A medida que el volumen de operaciones crece y los procesos se diversifican, la compañía enfrenta varios desafíos operativos y administrativos.

Actualmente, ChibchaWeb carece de una solución unificada que automatice y centralice sus procesos comerciales, operativos y de atención al cliente. Esto genera ineficiencias en la gestión de perfiles, pagos, comisiones, atención técnica y comunicación con registradores externos. La falta de un sistema integral limita la capacidad de escalar operaciones, introducir nuevos métodos de pago, y garantizar trazabilidad y control sobre los procesos internos.

\subsection{Descripción general del sistema a desarrollar}
El sistema a desarrollar para ChibchaWeb es una plataforma integral de gestión de servicios de hospedaje web, diseñada para automatizar, centralizar y optimizar los procesos operativos clave de la compañía. Esta solución tecnológica permitirá gestionar clientes, empleados, distribuidores, dominios, pagos y soporte técnico de manera eficiente, segura y escalable.

El sistema contará con módulos para la creación y administración de perfiles de clientes y empleados, y almacenará la información en una base de datos relacional optimizada incluyendo la validación de información personal y financiera garantizando la integridad evitando ciclos en esta; buscar y adquirir dominios mediante un carrito; adquirir paquetes de hosting.\\

El sistema contará con roles diferenciados por tipo de cuenta; búsquedas y consultas detalladas de todos los actores del sistema, así como herramientas para la generación de archivos de transferencia electrónica y seguimiento de actividades internas del proyecto; gestión de solicitudes de dominios con generación de archivos XML para su envío a registradores externos; Cálculo automático de comisiones a distribuidores según su categoría (BÁSICO o PREMIUM); y un módulo de soporte técnico basado en tickets que garantiza el seguimiento y resolución estructurada de problemas reportados.\\

La solución incluirá una interfaz web en React, backend en FastAPI y base de datos en MySQL. Se aplicarán buenas prácticas de diseño para mantener una l ógica de negocio clara y una estructura coherente en el modelo de datos.

\subsection{Objetivos de la solución}

\subsubsection{Objetivo general}
\begin{itemize}
	\item Desarrollar un sistema web que permita registrar, gestionar y monitorear cuentas de usuarios, adquisición de dominios y facturación de manera eficiente y centralizada.

\end{itemize}

\subsubsection{Objetivos específicos}
\begin{itemize}
\item Diseñar una base de datos relacional normalizada y sin ciclos innecesarios para soportar la lógica del sistema.
\item Desarrollar una interfaz web funcional y responsiva utilizando React y Bootstrap para facilitar la interacción del usuario con el sistema.
\item Construir un backend en FastAPI que gestione la lógica del negocio, la autenticación de usuarios y la comunicación con la base de datos.
\end{itemize}

\subsection{Stakeholders}
%configutación tablas
\pgfplotstableset{
    % Estilo general
    every head row/.style={
        before row=\hline\rowcolor{cafet}\bfseries,
        after row=\hline,
    },
    every odd row/.style={before row=\rowcolor{gray!10}},
    columns/Stakeholder/.style={string type, column type=p{4cm}},
    columns/Descripción/.style={string type,column type=p{12cm},},
    col sep=semicolon,
}

\begin{center}
    \begin{table}[H]
    \pgfplotstabletypeset[
    col sep=semicolon,
    header=true,
    ]{generalidades/stakeholders.dat}
    \caption{Stakeholders}
    \end{table}
\end{center}


\newpage
\section{Descripción y justificación de la metodología}
Para el desarrollo del sistema de información de ChibchaWeb, se ha seleccionado la metodología ágil Scrum. Esta elección se basa en las características del proyecto y las necesidades del cliente, que requieren flexibilidad, entregas parciales funcionales y un involucramiento constante de los interesados en el proceso.

\subsection{Scrum}
Scrum es un marco ágil de desarrollo iterativo e incremental que permite entregar valor de manera temprana y continua a través de sprints (iteraciones de tiempo fijo, típicamente de 1 a 4 semanas).\\
En Scrum, el trabajo se organiza en:

\begin{itemize}
    \item {Product Backlog: Lista priorizada de funcionalidades y requerimientos del sistema.}
    \item {Sprint Backlog: Conjunto de tareas seleccionadas para el Sprint en curso.}
    \item {Sprint Review y Sprint Retrospective: Evaluaciones al finalizar cada sprint, enfocadas en la mejora continua.}
\end{itemize}

\subsubsection{Roles en Scrum aplicados al proyecto}

\begin{itemize}
    \item {Product Owner: Representa los intereses de ChibchaWeb. Define y prioriza el Product Backlog.}
    \item {Scrum Master: Facilita el marco de trabajo Scrum, elimina impedimentos y promueve buenas prácticas.}
    \item {Development Team: Grupo autoorganizado de desarrolladores encargados de diseñar, codificar, probar e implementar el sistema.}
\end{itemize}

\subsection{Justificación de su uso en este proyecto}

\begin{itemize}
    \item {Adaptabilidad a cambios de requisitos: Dado que el negocio de hospedaje web evoluciona rápidamente (nuevas formas de pago, expansión internacional), Scrum permite adaptarse con agilidad.}
    \item {Entrega continua de valor: Con cada sprint, se entregan versiones funcionales del sistema que pueden ser probadas por el cliente.}
    \item {Comunicación constante con el cliente: La retroalimentación frecuente mejora la alineación entre lo que se construye y lo que el cliente realmente necesita.}
    \item {Reducción del riesgo: La entrega incremental permite detectar errores o desviaciones de manera temprana.}
    \item {Fomento de la colaboración y mejora continua: Las reuniones retrospectivas y la autoorganización del equipo fortalecen la eficiencia del proceso.}
\end{itemize}


\newpage
\section{Especificación y Recabación de Requerimientos Funcionales}
La presente sección describe los requerimientos funcionales necesarios para el desarrollo de la aplicación de software para ChibchaWeb, con el fin de automatizar y optimizar los procesos de gestión de clientes, pagos, dominios, distribuidores y soporte técnico.

\pgfplotstableset{
    % Estilo general
    every head row/.style={
        before row=\hline\rowcolor{cafet}\bfseries,
        after row=\hline,
    },
    every odd row/.style={before row=\rowcolor{gray!10}},
    columns/Código/.style={string type, column type=p{1.5cm}},
    columns/Requerimiento Funcional/.style={string type, column type=p{14.5cm}},
    col sep=semicolon,
}

\begin{center}
    \begin{table}[H]
    \pgfplotstabletypeset[
    col sep=semicolon,
    header=true,
    ]{especificacion/espec.dat}
    \caption{Requerimientos Funcionales, parte 1.}
    \end{table}
\end{center}

\begin{center}
    \begin{table}[H]
    \pgfplotstabletypeset[
    col sep=semicolon,
    header=true,
    ]{especificacion/espec2.dat}
    \caption{Requerimientos Funcionales, parte 2.}
    \end{table}
\end{center}


\section{Motivadores Arquitecturales}
\subsection{Motivadores de Negocio}
Los motivadores de negocio que se mostrarán a continuación definen las razones por las cuales ChibchaWeb requiere una nueva solución de software. Entre las cuales se encuentran:

\subsubsection{Expansión internacional del negocio}
ChibchaWeb, actualmente con clientes en Colombia y países vecinos, proyecta una expansión hacia mercados internacionales como África. Esto requiere una plataforma escalable y flexible que soporte la gestión de clientes, dominios y facturación en diferentes regiones.

\subsubsection{Automatización de procesos clave}
El sistema busca automatizar la creación, edición y eliminación de perfiles de clientes, empleados y distribuidores, reduciendo la carga operativa y los errores humanos, y facilitando la trazabilidad de operaciones críticas.

\subsubsection{Mejora en la atención al cliente y soporte técnic}

Al permitir el registro y seguimiento de tickets, el sistema mejorará el proceso de resolución de problemas, permitiendo al equipo de soporte priorizar incidencias y mejorar los tiempos de respuesta, alineándose con los niveles de servicio ofrecidos.

\subsubsection{Gestión y segmentación de distribuidores}
Con la existencia de distribuidores BÁSICOS y PREMIUM, el sistema debe permitir su categorización, cálculo de comisiones correspondientes (10% o 15%), y la generación de reportes financieros para pagos electrónicos, fomentando relaciones sólidas con aliados estratégicos.

\subsubsection{Flexibilidad en métodos de pago}
Actualmente se aceptan tarjetas de crédito, pero se planea incluir otros métodos como PSE. Esto requiere una arquitectura extensible que permita integrar fácilmente nuevos métodos de pago en el futuro sin afectar la operatividad actual.

\subsection{Restricciones de Tecnología}
El desarrollo de la plataforma para ChibchaWeb estará sujeto a las siguientes restricciones tecnológicas, derivadas de decisiones de arquitectura, compatibilidad, mantenibilidad y recursos disponibles:

Tecnologías definidas por el equipo de desarrollo

 \subsubsection{Frontend}

Se utilizará React.js como framework principal para la construcción de la interfaz de usuario, debido a su eficiencia, modularidad y gran comunidad de soporte.

\subsubsection{Backend}

Se empleará FastAPI con Python, por su rendimiento, facilidad de definición de servicios RESTful, y compatibilidad con documentación automática.

\subsubsection{Base de datos}

Se usará MySQL, una base de datos relacional robusta y ampliamente soportada, ideal para mantener la integridad referencial y relaciones entre entidades complejas (clientes, empleados, dominios, etc.).

\subsubsection{Control de versiones y colaboración}

El proyecto estará versionado mediante GitHub, lo cual permitirá control de cambios, trabajo colaborativo y despliegue continuo si se requiere.

\subsection{Restricciones de Negocio}
El desarrollo y funcionamiento del sistema estarán sujetos a diversas restricciones impuestas por las necesidades, normativas y condiciones del entorno comercial del proyecto. Estas restricciones afectan la manera en que se diseña y entrega la solución, e incluyen:

\subsubsection{Políticas de manejo de cuentas }

Solo se permitirán 4 tipos de cuentas: cliente, administrador, empleado y distribuidor, cada una con permisos y accesos claramente definidos. No se admitirán roles personalizados fuera de estas categorías durante la fase inicial del sistema.

\subsubsection{Manejo de dominios basado en disponibilidad externa }

La compra de dominios dependerá de su disponibilidad en registros externos. Si el dominio ya está tomado, no podrá ser adquirido ni facturado, lo cual impone una dependencia con proveedores externos y puede requerir lógica de validación en tiempo real.

\subsection{Atributos de Calidad}

Durante el desarrollo del sistema ChibchaWeb, se identificaron ciertos atributos de calidad que orientaron decisiones técnicas y de arquitectura. A continuación, se describen los principales, en función de la experiencia práctica con el sistema ya desplegado y en funcionamiento.

\subsubsection{Disponibilidad}

El sistema actualmente está desplegado mediante Railway, por lo tanto, su disponibilidad depende del servicio que ofrece esta plataforma. No se cuenta con una infraestructura propia, pero la integración entre base de datos, backend y frontend se mantiene activa y estable mientras Railway esté operativo. Esta disponibilidad es suficiente para los fines actuales, aunque si en un futuro fuera necesario se podría considerar una solución con mayor control sobre el entorno de producción.

\subsubsection{Seguridad}

Se implementó un sistema de autenticación por correo electrónico, en el cual los usuarios reciben un código de verificación. Solo quienes completen este proceso pueden acceder a las funciones reservadas para clientes. Esta medida ha resultado efectiva para restringir accesos y proteger la información del usuario, y se aplican validaciones en los formularios.


\subsubsection{Rendimiento}

La aplicación responde de forma fluida gracias a la conexión directa con las APIs del backend. No se ven demoras significativas al realizar operaciones como búsquedas, registros, etc.


\subsubsection{Escalabilidad}

Actualmente el sistema permite operaciones como búsqueda de dominios y la “compra” como tal, lo cual es una base para un entorno escalable. Aunque todavía no se ofrecen servicios de hosting reales desde la plataforma, el sistema está preparado para incluir esa funcionalidad más adelante, como tal el nucleo del sistema nos permite esto


\subsubsection{Mantenibilidad}

El código está organizado de forma clara, permitiendo modificar o añadir funcionalidades sin afectar lo ya implementado. Se han realizado ajustes menores y correcciones sin mayores inconvenientes, lo que demuestra que el sistema es mantenible. Además, se pueden actualizar componentes del frontend o backend de manera independiente.


\subsubsection{Usabilidad}

Desde la perspectiva del cliente, la interfaz es bastante intuitiva. Las funciones están bien distribuidas y no es necesario tener conocimientos técnicos para navegar o realizar operaciones básicas. En el caso del personal de soporte o empleados, hay más funciones que requieren cierta formación, como el manejo de tickets, pero están diseñadas de forma clara.


\subsubsection{Interoperabilidad}

El sistema se comunica con servicios externos mediante peticiones web, especialmente para consultar la disponibilidad de dominios. Para ello, se implementó un proceso de scraping que permite validar si un dominio está disponible o no.


\subsubsection{Recabación y justificació de Requerimientos No Funcionales}
Durante la planificación y construcción del sistema ChibchaWeb, se identificaron diversos requerimientos no funcionales a partir de las necesidades específicas del proyecto, la experiencia con despliegues reales, y las condiciones técnicas impuestas por el entorno en el que opera la aplicación. Estos requerimientos no son visibles directamente para el usuario final, pero son determinantes para que el sistema sea confiable, seguro y sostenible en el tiempo.

\begin{itemize}
    \item {Las operaciones realizadas por el usuario (como la búsqueda de dominios, la carga de formularios o la visualización de datos) no es de un tiempo significativo o excesivo, esto mejora la experiencia del usuario, esto gracias al uso de las Apis }
    \item {Aunque actualmente se encuentra todo en el servicio de railway para el despligue, fuera de eso el sistema es funcional y estable. }
    \item {Solo los usuarios verificados deben tener acceso a funcionalidades sensibles del sistema. Además, los datos sensibles  como contraseña están protegidos  }
    \item {El sistema esta preparado para soportar un aumento progresivo de usuarios y funcionalidades sin comprometer el rendimiento, aunque recalcar que al menos en esta fase la disponibilidad no depende tanto de nosotros, sino del servicio de railway }
    \item {En diferentes momentos del proyecto se han realizado cambios o mejoras con relativa facilidad, lo que indica que se ha logrado un diseño mantenible. }
    \item {Ya se implementó una primera forma de integración con servicios de disponibilidad de dominios usando scraping. Esta capacidad de interactuar con sistemas externos es la clave para estas funcionalidades. }
\end{itemize}

\subsubsection{Escenarios de Calidad}
Los escenarios de calidad permiten evaluar cómo se comporta el sistema ante situaciones específicas que afectan atributos no funcionales como el rendimiento, la seguridad o la disponibilidad. A continuación se presentan algunos escenarios relevantes que reflejan condiciones reales o esperadas para la plataforma ChibchaWeb:

\subsubsection*{Consulta masiva de dominios disponibles}

\begin{itemize}
\item \textbf{Contexto:} Un distribuidor accede a la plataforma y realiza varias búsquedas consecutivas para verificar la disponibilidad de dominios para sus clientes.

\item \textbf{Estimulación:} Se ejecutan 10 a 15 búsquedas en un corto periodo de tiempo.

\item \textbf{Respuesta esperada:} El sistema podra responder a estas consultas en un tiempo no tan significativo, sin errores ni bloqueos.
\end{itemize}

\subsubsection*{Intento de acceso sin verificación de correo}

\begin{itemize}
\item \textbf{Contexto:} Un usuario se registra pero no verifica su correo electrónico.

\item \textbf{Estimulación:} Intenta acceder a las funcionalidades de cliente como compra de dominio o visualización del perfil.

\item \textbf{Respuesta esperada:} El sistema bloquea el acceso y solicita la verificación de correo antes de continuar.
\end{itemize}

 \subsubsection*{Actualización del backend sin detener el sistema completo}

\begin{itemize}
\item \textbf{Contexto:} Se implementa una mejora en una API del backend.

\item \textbf{Estimulación:} Se despliega el nuevo backend mientras la base de datos y el frontend siguen activos.

\item \textbf{Respuesta esperada:} El sistema sigue funcionando normalmente o se interrumpe solo por un corto tiempo. No se generan inconsistencias.
\end{itemize}

\subsubsection*{Interacción de un cliente no técnico}

\begin{itemize}
\item \textbf{Contexto:} Un cliente nuevo accede al sistema para registrar su primer dominio.

\item \textbf{Estimulación:} Usa la interfaz para buscar un dominio, registrar sus datos y simular un pago.

\item \textbf{Respuesta esperada:} El cliente puede completar la operación sin necesidad de asistencia. Los pasos son claros y comprensibles.
\end{itemize}

\subsubsection*{Caída temporal del servicio Railway}

\begin{itemize}
\item \textbf{Contexto:} Railway presenta una interrupción temporal.

\item \textbf{Estimulación:} Usuarios intentan ingresar al sistema durante el tiempo en que el servicio está inactivo.

\item \textbf{Respuesta esperada:} El sistema no responde (fuera de servicio), pero se recupera automáticamente cuando Railway vuelve a estar disponible, sin pérdida de información.
\end{itemize}


\newpage
\section{Contexto}
\subsection{Escenarios operacionales}
Estos Escenarios operacionales representan las interacciones comunes o típicas que pasan entre los diferentes actores de la plataforma de ChibchaWeb, Cada escenario describe como un usuario o actor logra el objetivo deseado utilizando el sistema en un contexto especifico

\subsubsection{Escenario 1: Registro y verificación de cuenta}

\begin{itemize}
\item Actor: Usuario

\item Descripción: Un usuario se registra en la plataforma ingresando sus datos personales. Luego, valida su cuenta a través de un token enviado por correo electrónico.

\item Resultado esperado: El usuario puede iniciar sesión y acceder a las funciones de su rol.
\end{itemize}

\subsubsection{Escenario 2: Compra de un paquete de hosting}

\begin{itemize}
\item Actor: Cliente

\item Descripción: El cliente agrega un método de pago, selecciona un paquete, realiza el pago y se activa el servicio.

\item Resultado esperado: El sistema registra la compra.
\end{itemize}

\subsubsection{Escenario 3: Registro de solicitud de dominio}

\begin{itemize}
\item Actor: Usuario (Cliente o Distribuidor)

\item Descripción: El usuario busca un dominio, valida su disponibilidad y envía una solicitud.

\item Resultado esperado: La solicitud es procesada y almacenada para generación de archivos XML.
\end{itemize}

\subsubsection{Escenario 4: Atención de ticket de soporte}

\begin{itemize}
\item Actor: Usuario

\item Descripción: El usuario crea un ticket describiendo su problema. Soporte lo atiende, lo puede escalar y responder. El administrador puede asignarlo.

\item Resultado esperado: se le da solución al ticket lo más pronto posible
\end{itemize}



\subsubsection{Escenario 5: Gestión administrativa de usuarios}

\begin{itemize}
\item Actor: Administrador

\item Descripción: Desde el panel de administración, el administrador puede gestionar precios, paquetes y consultar

\item Resultado esperado: Los cambios se reflejan en la base de datos y actualizan el aplicativo en tiempo real
\end{itemize}



\subsubsection{Escenario 6: Cálculo y pago de comisiones}

\begin{itemize}
\item Actor: Administrador

\item Descripción: El sistema calcula las comisiones para los distribuidores según su categoría básico o premium (recordar que esto es un descuento)

\item Resultado esperado: El administrador tiene listos los valores y la documentación de soporte.
\end{itemize}

\subsubsection{Escenario 7: Consulta de perfil personal}

\begin{itemize}
\item Actor: Usuario

\item Descripción: El usuario inicia sesión, accede a su perfil y visualiza información relevante como nombre, correo, y acceso a sus métodos de pago o registrar uno nuevo

\item Resultado esperado: El usuario tiene acceso a un panel completo sobre sus datos de forma clara
\end{itemize}

\subsubsection{Escenario 8: Transferencia de dominio vía correo}

\begin{itemize}
\item Actor: Usuario

\item Descripción: Un usuario puede transferir un dominio en su propiedad al correo de otro usuario.

\item Resultado esperado: La transferencia queda registrada y se notifica a los usuarios involucrados
\end{itemize}

\subsection{Casos de Uso: Descripción y Modelo}

Los casos de uso del sistema ChibchaWeb describen las funcionalidades clave que pueden ser ejecutadas por los distintos actores, permitiendo alcanzar objetivos específicos mediante la interacción con la aplicativo. Este modelo facilita la comprensión de los límites del sistema, las responsabilidades de cada actor y las relaciones entre funcionalidades.\\

Actores Del Sistema:
\begin{itemize}
    \item Usuario: Actor general del que heredan todos los demás.
    \item Cliente: Usuario registrado que puede adquirir servicios y registrar dominios.
    \item Distribuidor: Usuario que gestiona múltiples dominios y cobra a sus propios clientes.
    \item Empleado: Usuario interno con funcione operativas.
    \item Soporte: Subtipo de empleado, encargado de atender tickets.
    \item Administrador: Usuario con permisos totales para gestionar perfiles, dominios y comisiones
\end{itemize}

\subsubsection{Resumen De Casos De Uso}

\pgfplotstableset{
    % Estilo general
    every head row/.style={
        before row=\hline\rowcolor{cafet}\bfseries,
        after row=\hline,
    },
    every odd row/.style={before row=\rowcolor{gray!10}},
    columns/Categoría/.style={string type, column type=p{4cm}},
    columns/Caso de Uso/.style={string type, column type=p{8cm}},
    columns/Actor(es)/.style={string type,column type=p{3cm},},
    col sep=semicolon,
}

\begin{center}
    \begin{table}[H]
    \pgfplotstabletypeset[
    col sep=semicolon,
    header=true,
    ]{contexto/rescas.dat}
    \caption{Resumen De Casos De Uso}
    \end{table}
\end{center}


\newpage
\section{Descripción y justificación de patrones de diseño adoptados y arquitectura de software definida.}
En el desarrollo de la plataforma ChibchaWeb se implementaron diversos patrones de diseño que permitieron mantener la claridad en la arquitectura, separación de responsabilidades y facilidad de mantenimiento del código.

\subsubsection{Patrón DTO (Data Transfer Object)}
Utilizado para la validación y transporte de datos entre el frontend y el backend. Se implementó mediante Pydantic en FastAPI, garantizando que la información recibida cumpla con el formato esperado antes de ser procesada, reduciendo así errores y mejorando la seguridad de entrada de datos.

\subsubsection{Patrón Active Record}
Aplicado en los modelos de datos mediante SQLAlchemy, permitiendo que cada entidad del sistema (Cliente, Dominio, Paquete, etc.) contenga tanto su representación como las operaciones CRUD asociadas, simplificando el acceso y manipulación de la base de datos.

\subsubsection{Patrón Controller}
Implementado a través de los routers de FastAPI, donde cada módulo del sistema cuenta con controladores que gestionan las solicitudes HTTP, orquestan la lógica de negocio y retornan las respuestas correspondientes.

\subsubsection{Encapsulación de lógica en servicios}
Funcionalidades específicas, como el generador de respuestas mediante inteligencia artificial, se desarrollaron como servicios independientes. En este componente se emplearon los patrones Strategy (para permitir distintas formas de generar respuestas según el contexto) y Adapter (para integrar proveedores externos de IA con la lógica interna del sistema).

\subsubsection{Patrón Template Method}
Los flujos de facturación y la generación de solicitudes XML para registradores externos se estructuraron siguiendo este patrón, definiendo la secuencia general de pasos y permitiendo que subclases concreten detalles específicos, lo que facilita cambios futuros sin modificar la estructura base del proceso.

En conjunto, la aplicación de estos patrones ha permitido desarrollar un sistema modular, escalable y adaptable a los cambios, manteniendo buenas prácticas de ingeniería de software y asegurando la coherencia entre los distintos componentes.


\newpage
\section{Puntos de Vista y Modelos Arquitecturales}
\subsection{Punto de Vista: Descripción del problema.}
\subsubsection{Punto de Vista: Contexto de la solución (Diagrama de Contexto)}
\subsection{Punto de Vista: Interacción}
\subsubsection{Diagrama de Casos de Uso}
\subsubsection{Diagramas de Secuencia}
\subsection{Punto de Vista: Desarrollo (Diagrama de Clases o cualquier otro)}
\subsection{Punto de Vista: Modelo Relacional de la BBDD.}


\newpage
\section{Firmas de aceptación con fecha}
Andrés José Acevedo Ardila
\begin{figure}[H]
    \centering
    \includegraphics[width=0.2\linewidth]{firmas/firma-andres.png}
\end{figure}

Brayan Steven Sánchez Casilima
\begin{figure}[H]
    \centering
    \includegraphics[width=0.2\linewidth]{firmas/firma-steveb.png}
\end{figure}

Christian Camilo Lancheros Sánchez
\begin{figure}[H]
    \centering
    \includegraphics[width=0.2\linewidth]{firmas/firma-cris.png}
\end{figure}

Juan Esteban Oviedo Sandoval
\begin{figure}[H]
    \centering
    \includegraphics[width=0.2\linewidth]{firmas/firma-esteban.png}
\end{figure}

Karen Tatiana Bravo Rodríguez
\begin{figure}[H]
    \centering
    \includegraphics[width=0.2\linewidth]{firmas/firma-karen.png}
\end{figure}

Luis Felipe Mayorga Tibaquicha
\begin{figure}[H]
    \centering
    \includegraphics[width=0.2\linewidth]{firmas/firma-felipe.png}
\end{figure}

Firmado el día:
07 de Agosto del año 2025


\newpage
\bibliographystyle{plain}
\bibliography{bibliografia/bibliografia}

\end{document}
